\chapter{Overview}
\label{ch:overview}

\Blackship{} is a FreeBSD jail orchestrator. It turns a TOML manifest and
optional Jailfiles into repeatable jail lifecycle operations: acquire a
userland, create or clone a root filesystem, configure isolation and resources,
attach networking, run hooks, supervise health, and reclaim every resource when
the jail stops.

The project sits between FreeBSD's native tools and a full cluster platform. It
does not replace \command{jail(8)}, ZFS, PF, rctl, cpuset, or rc; it coordinates
them and records enough state to recover from partial failure.

\section{Who this is for}

Blackship is designed for:

\begin{itemize}
  \item FreeBSD administrators who want declarative, dependency-aware jails;
  \item application teams running several cooperating services on one host;
  \item CI systems that need fast, disposable FreeBSD environments;
  \item management interfaces that need stable JSON rather than scraped text;
  \item developers who want a Rust library for ephemeral jail workloads; and
  \item operators who value native FreeBSD facilities over a parallel runtime.
\end{itemize}

It is particularly effective on ZFS hosts. A bootstrapped release can be
snapshotted once and cloned for each jail, making creation fast and
space-efficient. The non-ZFS path remains supported, but uses full directory
copies and cannot offer snapshots, clones, or ZFS send/receive.

\section{What it is not}

Blackship is not a Linux container compatibility layer, a virtual-machine
manager, a multi-host scheduler, or a Kubernetes implementation. A jail shares
the FreeBSD kernel with its host. Workloads that need a different kernel,
hardware emulation, live migration, or hostile multi-tenant isolation should be
evaluated against bhyve or another VM boundary.

Blackship also does not invent a package or service manager inside the jail.
Jailfiles run ordinary FreeBSD commands, and managed jails normally boot through
\file{/etc/rc} and stop through \file{/etc/rc.shutdown}.

\section{Major capabilities}

\begin{description}[style=nextline]
  \item[Declarative manifests]
    Describe jails, dependencies, hooks, VNET, health checks, resource limits,
    QoS phases, draining, and storage in \file{blackship.toml}.
  \item[Dependency-aware lifecycle]
    Start prerequisites first and stop dependents first. Select a jail by name,
    unambiguous prefix, tag, or \command{ALL}.
  \item[Release and image management]
    Bootstrap FreeBSD releases from distribution sets or pkgbase, build Jailfiles
    with a ZFS layer cache, and commit configured roots as reusable releases.
  \item[Native networking]
    Create VNET jails using either \command{if\_bridge}/epair or
    \command{ng\_bridge}/\command{ng\_eiface}, allocate addresses, configure DNS,
    and publish services through PF anchors.
  \item[Operational control]
    Apply live-safe changes with EVA, supervise health, enforce rctl and cpuset
    policy, drain connections, inspect audit history, and garbage-collect
    orphaned resources.
  \item[Integration surfaces]
    Use versioned JSON envelopes, Armada multi-file orchestration, a Rust library,
    shell completion, or the optional MCP server.
\end{description}

\section{A typical workflow}

The normal path is deliberately short:

\begin{enumerate}
  \item install Blackship on a supported FreeBSD host;
  \item bootstrap a base release;
  \item define one or more jails in \file{blackship.toml};
  \item create a named network when VNET is required;
  \item validate and preview the plan;
  \item start the target or the entire dependency graph; and
  \item observe, update, snapshot, export, or stop it using the same manifest.
\end{enumerate}

\begin{shipnote}[One source of truth]
The manifest describes desired jail configuration. Runtime stores beneath the
data directory record what Blackship actually created. Do not edit those stores
by hand; use the CLI, or reconcile manual host changes with \command{eva},
\command{cleanup}, or \command{gc}.
\end{shipnote}

\section{Documentation conventions}

The manual uses \emph{service name} for the short name in the manifest and
\emph{full jail name} for the project-prefixed kernel name. The word
\emph{release} means a bootstrapped FreeBSD userland or a committed compatible
image, not the host kernel release. A \emph{named network} is a persisted
Blackship record that can be recreated after reboot; it is not merely a bridge
interface that happens to exist.
